\chapter{Conclusion}
\label{chap:conclusion}

\section{Retour sur les objectifs SMART}

Les cinq objectifs formulés en C1
(\filepath{docs/architecture/etude_faisabilite.md} §7,
chapitre~\ref{chap:introduction}) sont atteints à la clôture du
programme (M0 à M4) :

\begin{enumerate}
  \item \textbf{Centraliser les données FluxPro, TransFlow et fichiers
    clients} -- base de travail unique, les 7 tables FluxPro et les
    sources TransFlow/fichiers clients intégrées (chapitre~\ref{chap:bloc2}).
  \item \textbf{Exposer une API REST documentée} -- Omega Data API,
    spécification OpenAPI auto-générée, autorisation par rôle
    (chapitre~\ref{chap:bloc2} §C12).
  \item \textbf{Mettre à disposition un tableau de bord} -- livré en
    deux temps, pas en une seule fois : vues SQL + notebook à la
    clôture du Bloc 3, puis Grafana (chapitre~\ref{chap:m4}) après
    qu'une relecture a établi que la première forme ne suffisait pas
    au sens littéral du référentiel. L'objectif est atteint, mais pas
    du premier coup -- à noter honnêtement plutôt que de présenter une
    trajectoire linéaire qui n'a pas eu lieu.
  \item \textbf{Réduire à zéro les traitements de données personnelles
    non répertoriés} -- registres RGPD tenus à chaque brique (staging,
    entrepôt, lake), audits automatisés plutôt que déclaratifs (voir
    le fil transverse du chapitre~\ref{chap:synthese-transverse}).
  \item \textbf{Concevoir une architecture de data lake} pour le flux
    IoT -- zones \texttt{raw}/\texttt{staging}/\texttt{curated}
    opérationnelles, ingestion batch et temps réel toutes deux vérifiées
    en conditions réelles (chapitre~\ref{chap:bloc4}).
\end{enumerate}

\section{Limites connues}

Consignées explicitement dans les comptes rendus de chaque jalon
(\filepath{docs/comptes_rendus/}) plutôt que découvertes a posteriori
pour les besoins de ce rapport :

\begin{itemize}
  \item \textbf{Rapprochement \texttt{clients.nom}~$\leftrightarrow$~\texttt{historique\_expeditions.client}}
    jamais garanti par une contrainte formelle -- vérifié cohérent sur
    le jeu de données actuel, pas structurellement sûr
    (chapitre~\ref{chap:bloc2}, point ouvert transmis dès M1).
  \item \textbf{Aucun ordonnanceur actif} pour les tâches planifiées
    (sauvegardes de l'entrepôt, purge RGPD du lake) -- scripts et
    procédures documentés, mais à exécuter manuellement ou via cron
    en conditions réelles (\filepath{docs/comptes_rendus/M2.md},
    \filepath{docs/comptes_rendus/M3.md}).
  \item \textbf{Stabilité de la clé de pseudonymisation entre
    exécutions} non garantie ni surveillée -- détaillé
    chapitre~\ref{chap:synthese-transverse}.
  \item \textbf{Prestataire SupervIT jamais mobilisé} -- le périmètre
    pédagogique du Bloc 4 n'en a pas eu besoin, écart
    budgétaire/organisationnel assumé plutôt qu'un manque non vu.
  \item \textbf{Observabilité partielle par choix} (pas d'agrégation
    centralisée des logs, pas de traçage distribué) et \textbf{pas
    d'orchestrateur de pipelines} -- les deux déjà justifiés comme des
    choix de sobriété délibérés plutôt que des manques non identifiés
    (chapitre~\ref{chap:m4}).
\end{itemize}

\section{Apprentissages}

Trois pratiques, imposées par l'expérience plutôt que décidées en
théorie au départ, se sont révélées transférables sur l'ensemble du
programme : \textbf{décider avant de coder} (un document de
séquencement par bloc, \filepath{docs/architecture/sequencement_bloc\{2,3,4\}.md},
écrit et daté avant la première ligne de code de chaque bloc) ;
\textbf{vérifier empiriquement plutôt que supposer} (le grain réel des
doublons mesuré plutôt qu'estimé, la cardinalité du rapprochement
livraison/expédition vérifiée à 100~\% avant de trancher le modèle de
données, les citations du référentiel de compétences relues au mot
près plutôt que travaillées de mémoire) ; et \textbf{la relecture
externe comme détecteur d'angles morts}, pas comme formalité -- la
restitution BI du jalon M4 et la faille de repli codé en dur sur la
clé de pseudonymisation (chapitre~\ref{chap:synthese-transverse}) ont
toutes deux été trouvées par un regard extérieur après qu'une première
analyse, faite seul, les avait jugées suffisantes ou écartées.

Le programme a par ailleurs été mené avec l'assistance d'outils
d'intelligence artificielle sous supervision humaine constante
(chapitre~\ref{chap:introduction}) : un apprentissage à part entière a
été de rester en mesure d'expliquer et de défendre chaque choix
technique du dépôt -- pas seulement de l'avoir accepté -- condition
nécessaire pour soutenir ce travail à l'oral.
